\section{Figures of stacked and fair trees}
\label{app:figures}


\begin{figure}[h!]
    \centering
    \scalebox{0.7}{
    \begin{tikzpicture}[tips=proper]
        \node[isosceles triangle,
            draw,thick,
            isosceles triangle apex angle=60,
            rotate=90,
            minimum size=3.5cm] (T1) at (0,0){};
        
        \node[isosceles triangle,
            draw,thick,
            isosceles triangle apex angle=60,
            rotate=90,
            minimum size=3cm,
            anchor=east] (T2L) at (T1.left corner){};

        \node[isosceles triangle,
            draw,thick,
            isosceles triangle apex angle=60,
            rotate=90,
            minimum size=3cm,
            anchor=east] (T2R) at (T1.right corner){};
    
        \node[isosceles triangle,
            draw,
            dashed,
            isosceles triangle apex angle=60,
            rotate=90,
            minimum size=11.52cm,
            anchor=east] (T5) at (T1.east){};
        
        \node[isosceles triangle,
            draw,
            thick,
            isosceles triangle apex angle=60,
            rotate=90,
            minimum size=3cm,
            anchor=east] (T3LL) at ([yshift=-0.00cm]T5.99){};
        \node[isosceles triangle,
            draw,
            thick,
            isosceles triangle apex angle=60,
            rotate=90,
            minimum size=3cm,
            anchor=east] (T3L) at ([xshift=3.5cm,yshift=-0.00cm]T5.99){};
        \node[isosceles triangle,
            draw,thick,
            isosceles triangle apex angle=60,
            rotate=90,
            minimum size=3cm,
            anchor=east] (T3R) at ([yshift=-0.00cm]T5.261){};    
        
        % dot dot dot
        \draw[] (0,-3) node [] {\LARGE $\cdots$};
        \draw[] (1.7,-7.6) node [] {\LARGE $\cdots$};
        \draw[] ([yshift=-0.5cm]T2L.west) node [rotate=90] {\LARGE $\cdots$};
        \draw[] ([yshift=-0.5cm]T2R.west) node [rotate=90] {\LARGE $\cdots$};

        % triangle labels
        \draw[] (T1.center) node [] {$T_1$};
        \draw[] (T2L.center) node [] {$T_2$};
        \draw[] (T2R.center) node [] {$T_2$};
        \draw[] (T3L.center) node [] {$T_k$};
        \draw[] (T3LL.center) node [] {$T_k$};
        \draw[] (T3R.center) node [] {$T_k$};

        % Depth label
        \draw[{Stealth[scale=1]}-{Stealth[scale=1]}] ([xshift=3.1cm]T1.east) to node[midway,fill=white!30,scale=1] {depth $d$} ([xshift=3cm]T1.west);
               
    \end{tikzpicture}
    }
  \captionsetup{width=.9\linewidth}

    \caption{Illustration of a stacked decision tree for a function $f^{\otimes k}$. The decision tree consists of $k$ depth-$d$ decision trees, $T_1,\ldots, T_k$, stacked on top of each other. For an input $X \in (\bn)^k$, the output $T(X)$ is computed sequentially, first by computing $T_1(X)$, then $T_2(X)$, and so on. The final output is $T(X)\coloneqq (T_1(X),\ldots, T_k(X))$. }
    \label{fig:stacked}
\end{figure}



\begin{figure}[h!]
    \centering
    \begin{tikzpicture}[tips=proper]
        \node[isosceles triangle,
            draw,
            thick,
            isosceles triangle apex angle=60,
            rotate=90,
            minimum size=6cm] (T1) at (0,0){};
        
        \draw[black,dashed] (T1.east) .. controls ([xshift=-0.1cm]T1.358) .. ([yshift=-0.5cm]T1.east) node[] (N1) {{}};
        \draw[black,dashed] ([yshift=-0.5cm]T1.east) .. controls ([xshift=0.4cm]T1.40) and (T1.350) .. (T1.center) node[] (N1) {{}};
        \draw[black,dashed] (T1.center) .. controls ([xshift=2.4cm]T1.110) .. (T1.west) node[] (N2) {{}};

        % different xi nodes
        \node[draw,circle,inner sep=0pt,fill=white] (x3) at (-0.45,-1.2) {\footnotesize $X_{8}^{(i)}$};
        \node[draw,circle,inner sep=0pt,fill=white] (x2) at (0,0) {\footnotesize $X_{3}^{(i)}$};
        \node[draw,circle,inner sep=0pt,fill=white] (x1) at (0,2) {\footnotesize $X_{22}^{(i)}$};
        
        % path label
        \draw[color=black] (T1.west) node [below,fill=white] {$\pi$};
        \node[draw,circle,fill=black,inner sep=1pt] (x) at (T1.west) {};

        % label to xi nodes
        % \node[fill=white] (label) at (-3,0.5) {\small $x^{(i)}$ queried $\le d$ times on path $\pi$};
        % \draw[-{Stealth[scale=1]}] (label) edge [out=0, in=200] (x1);
        % \draw[-{Stealth[scale=1]}] (label) edge [out=0, in=140] (x2);
        % \draw[-{Stealth[scale=1]}] (label) edge [out=0, in=130] (x3);
        \draw[] (-3.7,1) node [fill=white,text width=26mm] {\small $X^{(i)}$ queried $\le d$ times on path $\pi$};
        \coordinate (p) at (-2.3,1);
        \draw[-{Stealth[scale=0.75]},dotted,thick] (p) to (x1);
        \draw[-{Stealth[scale=0.75]},dotted,thick] (p) to (x2);
        \draw[-{Stealth[scale=0.75]},dotted,thick] (p) to (x3);

    \end{tikzpicture}
  \captionsetup{width=.9\linewidth}
    \caption{Illustration of a fair decision tree. For every block $i\in [k]$ and path $\pi$, the input block $X^{(i)}$ is queried at most $d$ times.} 
    \label{fig:fair}
\end{figure}



\section{Proof of~\Cref{thm:hardcore theorem for DTs}}
\label{appendix:hardcore}
Let $\mathcal{H}$ denote the set of measures of density $\delta/2$ with respect to $\mu$ and let $\mathcal{T}$ denote the set of decision trees $T$ whose expected depth with respect to $\mu$ is at most $d$. Suppose for contradiction that there \textit{does not} exist an $H\in\mathcal{H}$ which is $(\gamma,d)$-hardcore. That is, for all $H\in\mathcal{H}$ there is a tree $T$ of expected depth at most $d$ satisfying
\begin{equation}
\label{eq:minimax}
\Ex_{\bx\sim\mu}[f(\bx)T(\bx)H(\bx)]>\gamma\Ex_{\bx\sim\mu}[H(\bx)]=\gamma\delta/2.
\end{equation}
We use the minimax theorem to switch the quantifiers in the above statement. Consider the payoff matrix $M$ whose rows are indexed by distributions from $\mathcal{H}$ and whose columns are indexed by algorithms from $\mathcal{T}$ and whose entries are given by $M_{H,T}\coloneqq \Ex_{\bx\sim\mu}[f(\bx)T(\bx)H(\bx)]$. This is the payoff matrix for the zero-sum game where the row player first chooses a row $H$ and the column player then chooses a column $T$ and the payoffs are determined by $M_{H,T}$. Note that once the first player's strategy is fixed, we can assume without loss of generality that the second player's strategy is deterministic. Therefore, the minimax theorem for zero-sum games yields
\begin{align*}
    \gamma\delta/2&<\min_{\rho\in\mu(\mathcal{H})}\max_{T\in\mathcal{T}}~ (\rho^{\top} M)_T\tag{\Cref{eq:minimax}}\\
    &=\max_{\tau\in\mu(\mathcal{T})}\min_{H\in\mathcal{H}}~(M\tau)_H\tag{minimax theorem}
\end{align*}
where $\mu(\cdot)$ denotes the set of distributions over a given set. Therefore, there is a fixed distribution $\tau$ over the set $\mathcal{T}$ such that for all $H\in\mathcal{H}$
\begin{equation}
\label{eq:hardcore}
\Ex_{\bT\sim\tau}\bigg[\Ex_{\bx\sim\mu}[f(\bx)\bT(\bx)H(\bx)]\bigg]>\gamma\delta/2.    
\end{equation}
This shows that
$$
\Prx_{\bx\sim\mu}\bigg[\Ex_{\bT\sim\tau}[\bT(\bx)]f(\bx)\ge \gamma\bigg]\ge 1-\delta/2.
$$
In particular, if instead $\Prx_{\bx\sim\mu}[\Ex_{\bT\sim\tau}[\bT(\bx)]f(\bx)<\gamma]\ge \delta/2$ then we can contradict \Cref{eq:hardcore} by constructing a $\delta/2$-density $H$ such that $H(x)\coloneqq\Pr_{\bx\sim\mu}[\bx=x]$ for a $\delta/2$-fraction of $x$ satisfying $\Ex_{\bT\sim\tau}[\bT(\bx)]f(\bx)< \gamma$. \Cref{eq:hardcore} shows that $\Ex_{\bT\sim\tau}[\bT(\bx)]$ has good correlation with $f$ for a large fraction of inputs. We obtain a single strategy from the distribution $\tau$ by sampling $\bT_1,\ldots,\bT_r\sim\tau$ for $r$ sufficiently large (chosen later) and defining $T^\star$ as $\bT^{\star}(x)\coloneqq \text{MAJ}(\bT_1(x),\ldots,\bT_r(x))$.
For every $x$ for which $\Ex_{\bT\sim\tau}[\bT(x)]f(x)\ge \gamma$, we have
$$
\Prx_{\bT_1,\ldots,\bT_r\sim\tau}\bigg[\text{MAJ}(\bT_1(x),\ldots,\bT_r(x))\neq f(x)\bigg]\le 2^{-\Omega(\gamma^2 r)}
$$
by a Chernoff bound. Choosing $r=\Theta(\log(1/\delta)/\gamma^2)$ ensures that the failure probability is at most $\delta/2$. The decision tree $\bT^{\star}$ satisfies $\Pr_{\bx\sim\mu}[\bT^\star(\bx)\neq f(\bx)]\le \delta$. The expected depth of $\bT^\star$ is less than
$$
r\cdot d=\Theta(d\log (1/\delta)/\gamma^2)<\overline{\mathrm{Depth}}^\mu(f,\delta)
$$
which is a contradiction.

\section{The lack of error reduction for distributional error}
\label{sec:no-boosting}
In \Cref{sec:challenges}, we showed how error reduction gave a simple proof of a strong direct sum theorem for \emph{randomized} query complexity. The specific statement needed in that proof is the following standard error reduction by repetition theorem.
\begin{fact}[Error reduction for $\overline{\mathrm{R}}$]
    For any function $f:\bits^n \to \bits$ and $\delta > 0$,
    \begin{equation*}
        \overline{\mathrm{R}}(f,\delta) \leq O(\log(\lfrac{1}{\delta}))\cdot \overline{\mathrm{R}}(f,1/4).
    \end{equation*}
\end{fact}
Here, we give a short proof that no error reduction holds in the distributional setting, even with substantially weaker parameters.
\begin{claim}
    \label{claim:no-boosting}
    For any $n \in \N$, let $\mu$ be the uniform distribution over $\bits^n$. There is a function $f:\bits^n \to \bits$ satisfying,
    \begin{equation*}
        \overline{\mathrm{Depth}}^{\mu}(f,1/4) = 0 \quad\quad\quad\text{and}\quad\quad\quad \overline{\mathrm{Depth}}^{\mu}(f,1/8) \geq \Omega(n).
    \end{equation*}
\end{claim}
Since any function on $n$ bits can be computed exactly using $n$, the $f$ in the above claim requires essentially the maximum number of queries to be computed to error $1/8$ despite requiring no queries to be computed to error $1/4$.
\begin{proof}
    We first define $f$: If $x_1 = 0$, then $f(x) = 0$. Otherwise, $f(x)$ is the parity of the remaining $n-1$ bits of $x$.

    Note that,
    \begin{equation*}
        \Pr_{\bx \sim \mu}[f(\bx) = 0] = \frac{1}{2} \cdot \paren*{\Pr[f(\bx) = 0 \mid \bx_1 = 0] + \Pr[f(\bx) = 0 \mid \bx_1 = 1]} = \frac{3}{4}.
    \end{equation*}
    Therefore, the $0$ query algorithm that simply outputs $0$ has an error of only $1/4$ on $f$. 
    
    It only remains to prove that $\overline{\mathrm{Depth}}^{\mu}(f,1/10) \geq \Omega(n)$. Consider any (potentially randomized) $T:\bits^n \to \bits$ and leaf $\ell$ of $T$ at depth strictly less than $n-1$. By \Cref{prop:par-query-all}
    \begin{equation*}
        \Pr_{\bx \sim \mu}[T(\bx) \neq f(\bx) \mid \bx \text{ reaches } \ell, \bx_1 = 1] = 1/2.
    \end{equation*}
    Let $p$ be the probability that $T(\bx)$ queries a leaf of depth at least $n-1$ given that $\bx_1 = 1$. The above allows us to conclude that
    \begin{equation*}
        \Pr_{\bx \sim \mu}[T(\bx) \neq f(\bx)] \geq \frac{1}{2} \cdot \Pr[T(\bx) \neq f(\bx)\mid \bx_1 = 1] \geq \frac{1}{4} \cdot p. 
    \end{equation*}
    Therefore, if $T$ has error at most $1/8$, then $p$ must be at least $1/2$, which shows that $\overline{\mathrm{Depth}}^{\mu}(f,1/8) \geq \frac{n-1}{4}$.
\end{proof}